\documentclass[10pt,a4paper]{amsart}
\usepackage[margin=19mm]{geometry}
\usepackage{amsmath,amssymb,mathtools}
\usepackage[scr=rsfs,cal=euler]{mathalfa}
\usepackage{xfrac}
\usepackage[unicode,hidelinks,bookmarks=false]{hyperref}
\newcommand{\ie}{i.e.\ }
\newcommand{\cf}{cf.\ }
\newcommand{\resp}{respectively\ }
\newcommand{\Subsection}{\subsection}
\providecommand{\nobreakdash}{\nobreak\hbox{-}}
\setlength{\emergencystretch}{2em}
\multlinegap0pt
\usepackage{bm}
\usepackage{bbm}
\usepackage{dsfont}
\usepackage{xspace}
\usepackage{cancel}
\usepackage{mathtools}
\usepackage{amsthm}
\makeatletter
\@ifundefined{theorem}{\newtheorem{theorem}{Theorem}}{}
\@ifundefined{lemma}{\newtheorem{lemma}[theorem]{Lemma}}{}
\makeatother
\def\dR{{\mathbb R}}
\def\lm{\lambda}
\def\p{\partial}
\def\sH{{\mathfrak H}}   \def\sI{{\mathfrak I}}
\def\sfM{{\mathsf M}}   \def\sfN{{\mathsf N}}   \def\sfO{{\mathsf O}}
\newcommand\Omg{\Omega}
\newcommand\cL{\mathcal L}
\newcommand\ov{\overline}
\newcommand\sm{\setminus}
\title[On shape optimization with large magnetic fields in two dime]{On shape optimization with large magnetic fields in two dimensions}
\author{Vladimir Lotoreichik, L\'eo Morin}
\date{Source: arXiv:2509.08412v1 (CC BY 4.0)}
\begin{document}
\maketitle

\noindent\emph{Source and licence.} On shape optimization with large magnetic fields in two dimensions. Version arXiv:2509.08412v1 (CC BY 4.0), distributed under the Creative Commons Attribution 4.0 International licence, as recorded in the arXiv licence metadata for this exact version and shown on the abstract page https://arxiv.org/abs/2509.08412v1. Full rights notice: licenses/arxiv-2509.08412-CC-BY-4.0.txt.\par\smallskip

\noindent\textit{Editorial scope.} Counted unit: the Lemma whose source label is \texttt{lem.lower.dirac} in the article (source0.tex lines 780--785, source label \texttt{lem.lower.dirac}), delivered together with the article's own complete original proof of it (source0.tex lines 786--810, one \texttt{proof} environment of 25 lines). No mathematics is written, repaired or re-derived: statement and proof are the article's own text line for line. Required context retained uncounted: eq:minmaxDirac (lines 763--765). Every definition, display and label of those lines is retained unchanged. Omitted from this package and not redistributed: 1--762, 766--779, 811--1122. Nothing inside a retained excerpt is omitted or altered: every retained body is delivered exactly as the article's lines are written, so no omission receipt of any kind accompanies this package. Citations: the retained excerpts carry no citation command at all, so no bibliography entry of the article is transcribed and no reference is redistributed. Reference adaptation: none was needed and none was made; every reference inside the delivered unit resolves inside it, so no unresolved reference or citation is produced. Printed slips kept literal: no printed word, symbol, hypothesis or number of the counted statement or of its proof was altered. Layout and class: the article's own class is \texttt{amsart} with packages including inputenc, fontenc, amsmath, amsfonts, amssymb, lmodern, amsthm, accents, setspace, xcolor, mathtools, esint, mathrsfs, hyperref; the wrapper is the lane's pinned \texttt{amsart} editorial wrapper at 10pt on A4, which restates the theorem-like environments the retained text needs and adds the article's own macro definitions that the retained text needs (\texttt{\textbackslash Omg}, \texttt{\textbackslash cL}, \texttt{\textbackslash dR}, \texttt{\textbackslash lm}, \texttt{\textbackslash ov}, \texttt{\textbackslash p}, \texttt{\textbackslash sH}, \texttt{\textbackslash sfM}, \texttt{\textbackslash sm}), with extra wrapper packages (bm, bbm, dsfont, xspace, cancel, mathtools). The delivered numbering is the unit's own. Assets: no retained span contains a figure environment, a graphics-inclusion command, a PDF-page inclusion, a table or a TikZ picture; no image file is copied into this package, the package declares no asset dependency and no image file is redistributed with it. Licence: version arXiv:2509.08412v1 of this article is distributed under the Creative Commons Attribution 4.0 International licence, as recorded in the arXiv licence metadata for this exact version and shown on the abstract page https://arxiv.org/abs/2509.08412v1. A full rights notice is delivered at licenses/arxiv-2509.08412-CC-BY-4.0.txt. Characters: every retained excerpt is pure ASCII, so the delivered engine, XeLaTeX with its default Latin Modern text font, reports no missing character.
\begin{equation}\label{eq:minmaxDirac}
	\lm_n^+(\Omg,B) = \inf_{\begin{smallmatrix} \cL\subset\sH_{{\bf A}}(\Omg)\\ \dim \cL = n\end{smallmatrix}}\sup_{u\in\cL\sm\{0\}} \frac{\|u\|^2_{L^2(\p\Omg)} + \sqrt{\|u\|^4_{L^2(\p\Omg)}+ 4\|u\|^2_{L^2(\Omg)} \|d_{\bf A}^\times u\|^2_{L^2(\Omg)}}}{2\|u\|^2_{L^2(\Omg)}}.
\end{equation}

\begin{lemma}\label{lem.lower.dirac}
	Let $\Omg\subset\dR^2$ be a bounded, simply-connected domain with $\mathscr C^2$ boundary. Then for any $B > 0$ we have
	\[
		\lm_1^+(\Omg,B) \ge \lm_1^+(\Omg,0) e^{-2B\varphi^\Omg_{\rm m}}.
	\]
\end{lemma}

\begin{proof}
	Throughout the proof we use the abbreviation $\sH(\Omg) := \sH_0(\Omg)$ for the auxiliary Hilbert space corresponding to the zero vector potential.	First, we observe that the
	bounded and boundedly invertible multiplication operator 
	$\sfM u := e^{-B\varphi^\Omg} u$ acting in $L^2(\Omg)$ is a bijection between $\sH_{\bf A}(\Omg)$ and $\sH_0(\Omg)$. Indeed, for any $u\in\sH_{\bf A}(\Omg)$ we have a decomposition $u = u_0 + u_1$, where $u_0\in H^1(\Omg)$ and $u_1\in L^2(\Omg)$ satisfies $d_{\bf A}^\times u_1 = 0$ and $u_1|_{\p\Omg}\in L^2(\p\Omg)$. By regularity of the torsion function we have $\sfM u_0 \in H^1(\Omg)$.
	Moreover, it holds that $u_1|_{\p\Omg} = \sfM u_1|_{\p\Omg}\in L^2(\p\Omg)$, where we used that torsion function vanishes on the boundary of the domain $\Omg$. Finally, we also notice
	that for any $w\in L^2(\Omg)$ with $d_{\bf A}^\times w\in L^2(\Omg)$ we have
	\begin{equation}\label{eq:identity}
		d_0^\times \sfM w = e^{-B\varphi^\Omg}\Big(-2i \p_{\ov z} w + i B(\p_1\varphi^\Omg) w  - B(\p_2\varphi^\Omg) w\Big) = e^{-B\varphi^\Omg} d_{\bf A}^\times w. 
	\end{equation}
	Thus, we conclude that $\sfM u =\sfM u_0 + \sfM u_1\in\sH(\Omg)$, since $\sfM u_0\in H^1(\Omg)$ and
	by the above identity $\sfM u_1\in\mathscr{H}_{0}^2(\Omg)$. Hence, we have established that $\sfM$ maps $\sH_{{\bf A}}(\Omg)$ into $\sH(\Omg)$. Analogously, one can check that $\sfM^{-1}$ maps $\sH(\Omg)$ into $\sH_{{\bf A}}(\Omg)$. Moreover, since $\sfM$ has no kernel, we conclude that $\sfM$ is a bijection between $\sH_{\bf A}(\Omg)$ and $\sH(\Omg)$.
	Hence, by the variational characterisation~\eqref{eq:minmaxDirac} we get
	\[
	\begin{aligned}
		\lm_1^+(\Omg,B) &= \inf_{u\in \sH_{{\bf A}}(\Omg)\sm\{0\}}
		\frac{\|u\|^2_{L^2(\p\Omg)} + \sqrt{\|u\|^4_{L^2(\p\Omg)}+ 4\|u\|^2_{L^2(\Omg)} \|d_{\bf A}^\times u\|^2_{L^2(\Omg)}}}{2\|u\|^2_{L^2(\Omg)}}\\
		&=\inf_{v\in \sH_0(\Omg)\sm\{0\}}
		\frac{\|v\|^2_{L^2(\p\Omg)} + \sqrt{\|v\|^4_{L^2(\p\Omg)}+ 4\|e^{B\varphi^\Omg}v\|^2_{L^2(\Omg)} \|e^{B\varphi^\Omg}d_0^\times v\|^2_{L^2(\Omg)}}}{2\|e^{B\varphi^\Omg} v\|^2_{L^2(\Omg)}}\\
		&\ge e^{-2B\varphi^\Omg_{\rm m}}\inf_{v\in \sH_0(\Omg)\sm\{0\}}
				\frac{\|v\|^2_{L^2(\p\Omg)} + \sqrt{\|v\|^4_{L^2(\p\Omg)}+ 4\|v\|^2_{L^2(\Omg)} \|d_0^\times v\|^2_{L^2(\Omg)}}}{2\|v\|^2_{L^2(\Omg)}}\\
				& = e^{-2B\varphi^\Omg_{\rm m}}\lm_1^+(\Omg,0),
		\end{aligned}
	\]
	where we used the properties of the mapping $\sfM$, in particular, we employed the identity~\eqref{eq:identity}.
\end{proof}

\end{document}
